\chapter{Shuffling Cards}

\begin{center}\itshape
How many random transpositions mix a deck of $n$ cards? Fourier analysis on $\mathfrak S_n$ turns the question into a sum over Young diagrams, and the character formula of Chapter~9 evaluates it: the answer is $\frac12n\ln n$, and the deck goes from ordered to
mixed abruptly, in a window of order $n$ swaps. This is the first proof of the cutoff phenomenon.
\end{center}

Persi Diaconis left school at about fourteen to travel with the magician Dai Vernon, described by Martin Gardner as ``a professional card shark, or a card mechanic, as they call it in the trade''. Years later, when a friend told him that the best book
on probability was William Feller's \emph{Introduction to Probability Theory and its Applications}, he tried it and found that he could not read it: ``I didn't know calculus, or at least not enough.'' So he took evening classes at City College in New York, paying
for them with magic shows. When he applied to graduate school, Gardner wrote to Frederick Mosteller at Harvard: ``This young student is one of the best card mechanics in the country.'' Mosteller answered: ``If he's willing to major in statistics, I can get him into
Harvard.'' Diaconis took his doctorate there in 1974 and went to Stanford.

In 1981, with Mehrdad Shahshahani, he solved a problem that a card mechanic could appreciate. Swap two random cards, again and again: when is the deck mixed? Tome~3, Chapter~8, proved that about $2n\ln n$ swaps suffice (Broder's marking argument) and
that $\frac12n\ln n-2n$ do not. Diaconis and Shahshahani closed the gap with representation theory: $\frac12n\ln n+cn$ swaps suffice, for a constant $c$. For a deck of $52$ cards, $\frac12n\ln n\approx103$. Their tool, a Fourier transform on the non-commutative
group $\mathfrak S_n$, is what Tome~2 (Chapter~4) and Tome~3 (Chapter~8) promised; their key formula is Theorem~\ref{thm:transposition}.

The answer has a striking shape. The distance to uniform stays near $1$ for $\frac12n\ln n-cn$ swaps, then falls to near $0$ within a few multiples of $n$ more (Figure~\ref{fig:cutoff}). David Aldous and Diaconis named this the \emph{cutoff phenomenon} in 1986; it
is the reason why ``seven riffle shuffles'' (Tome~3, Chapter~8) is a meaningful answer.

Throughout, $G$ is a finite group, and $\widehat G$ is a set of representatives of its irreducible representations $\rho:G\to\mathrm{GL}(V_\rho)$, with $d_\rho=\dim V_\rho$ and character $\chi_\rho$. Each $V_\rho$ carries an inner product for which $\rho(g)$ is unitary (average any inner product
over $G$, as in the proof of Maschke's theorem), and we use orthonormal bases, so that $\rho(g^{-1})=\rho(g)^*$, the conjugate transpose. We use the notation of Tome~3, Chapter~8: $\|\mu-\nu\|=\frac12\sum_g|\mu(g)-\nu(g)|$ is the total variation distance, and $d(t)$ the
distance to stationarity after $t$ steps, from the worst starting point.

\section{Fourier analysis on a finite group}

\begin{definition}\label{def:fouriergroup}
The \emph{Fourier transform} of a function $f:G\to\C$ at $\rho\in\widehat G$ is the matrix $\hat f(\rho)=\sum_{g\in G}f(g)\rho(g)$. The \emph{convolution} of $f_1,f_2$ is $(f_1*f_2)(g)=\sum_{h\in G}f_1(h)f_2(h^{-1}g)$.
\end{definition}

In other words, $\hat f(\rho)$ is the image of $\sum f(g)g\in\C[G]$ in the factor $\operatorname{End}(V_\rho)$ of Wedderburn's isomorphism (Corollary~\ref{cor:wedderburnCG}), and convolution is the product of $\C[G]$. So $\widehat{f_1*f_2}(\rho)=\hat f_1(\rho)\hat f_2(\rho)$, as one can also
check directly: $\sum_g\sum_hf_1(h)f_2(h^{-1}g)\rho(h)\rho(h^{-1}g)$ is $\hat f_1(\rho)\hat f_2(\rho)$ after the substitution $g=hx$. If $X$ and $Y$ are independent random elements of $G$ with laws $\mu$ and $\nu$, the law of $XY$ is $\mu*\nu$: $P(XY=g)=\sum_hP(X=h)P(Y=h^{-1}g)$.
For $G$ abelian, $\widehat G$ is the dual group, the matrices are $1\times1$, and this is the Fourier analysis of Tome~2, Theorem~11.14, up to a complex conjugation.

\begin{theorem}[Fourier inversion and Plancherel]\label{thm:plancherelG}
For $f,f':G\to\C$ and $g\in G$,
\[
f(g)=\frac1{|G|}\sum_{\rho\in\widehat G}d_\rho\operatorname{tr}\bigl(\rho(g^{-1})\hat f(\rho)\bigr),\qquad
\sum_{g\in G}f(g)\overline{f'(g)}=\frac1{|G|}\sum_{\rho\in\widehat G}d_\rho\operatorname{tr}\bigl(\hat f(\rho)\hat f'(\rho)^*\bigr).
\]
\end{theorem}

\begin{proof}
\emph{The regular character.} $G$ acts on $\C[G]$ by left multiplication, permuting the basis; $g\ne e$ fixes no basis element, so the character of this \emph{regular representation} is $\chi_{\mathrm{reg}}(g)=|G|$ if $g=e$ and $0$ otherwise. The multiplicity of
$\rho$ in it is $\langle\chi_{\mathrm{reg}},\chi_\rho\rangle=\frac1{|G|}|G|\overline{\chi_\rho(e)}=d_\rho$ (Tome~2, Theorem~11.8), so $\chi_{\mathrm{reg}}=\sum_\rho d_\rho\chi_\rho$, that is,
\[
\sum_{\rho}d_\rho\chi_\rho(x)=|G|\,\mathbf 1_{x=e}.
\]

\emph{Inversion.} Both sides are linear in $f$, so it suffices to take $f=\mathbf 1_h$, with $\hat f(\rho)=\rho(h)$. The right side is $\frac1{|G|}\sum_\rho d_\rho\chi_\rho(g^{-1}h)$, which is $1$ if $g=h$ and $0$ otherwise, by the identity above.

\emph{Plancherel.} Apply inversion to $f$, multiply by $\overline{f'(g)}$ and sum over $g$:
\[
\sum_gf(g)\overline{f'(g)}=\frac1{|G|}\sum_\rho d_\rho\operatorname{tr}\Bigl(\Bigl(\sum_g\overline{f'(g)}\rho(g)^*\Bigr)\hat f(\rho)\Bigr)=\frac1{|G|}\sum_\rho d_\rho\operatorname{tr}\bigl(\hat f'(\rho)^*\hat f(\rho)\bigr),
\]
using $\rho(g^{-1})=\rho(g)^*$; and $\operatorname{tr}(AB)=\operatorname{tr}(BA)$.
\end{proof}

\begin{proposition}\label{prop:classfourier}
If $f$ is constant on conjugacy classes, then $\hat f(\rho)$ is the scalar matrix $\frac1{d_\rho}\bigl(\sum_gf(g)\chi_\rho(g)\bigr)I$.
\end{proposition}

\begin{proof}
For $x\in G$, substituting $g\mapsto x^{-1}gx$,
\[
\rho(x)\hat f(\rho)\rho(x)^{-1}=\sum_gf(g)\rho(xgx^{-1})=\hat f(\rho).
\] So $\hat f(\rho)$ commutes with every $\rho(x)$, and is a scalar by Schur's lemma; taking traces identifies the scalar.
\end{proof}

\section{The upper bound lemma}

A random walk on $G$ driven by a probability $Q$ moves from $x$ to $\tau x$, where $\tau$ is drawn from $Q$ independently at each step. Started at $e$, its law after $t$ steps is the $t$-fold convolution $Q^{*t}$, with Fourier transform $\hat Q(\rho)^t$. Started at $x$, its law is the
translate $g\mapsto Q^{*t}(gx^{-1})$, at the same distance from the uniform law $U$; so $d(t)=\|Q^{*t}-U\|$. Note that $\hat U(\rho)=0$ for $\rho\ne1$: the sum $\sum_g\rho(g)$ commutes with $\rho(G)$, so it is a scalar, of trace $|G|\langle\chi_\rho,1\rangle=0$.

\begin{lemma}[Diaconis--Shahshahani upper bound lemma]\label{lem:upperbound}
For a probability $\mu$ on $G$,
\[
\|\mu-U\|^2\le\frac14\sum_{\rho\ne1}d_\rho\operatorname{tr}\bigl(\hat\mu(\rho)\hat\mu(\rho)^*\bigr),
\]
the sum running over the non-trivial irreducible representations.
\end{lemma}

\begin{proof}
By Cauchy--Schwarz, $\bigl(\sum_g|\mu(g)-U(g)|\bigr)^2\le|G|\sum_g|\mu(g)-U(g)|^2$. By Plancherel, applied to $f=f'=\mu-U$, the last sum is $\frac1{|G|}\sum_\rho d_\rho\operatorname{tr}(\hat f(\rho)\hat f(\rho)^*)$. Here $\hat f(1)=1-1=0$, and $\hat f(\rho)=\hat\mu(\rho)$ for $\rho\ne1$.
Divide by $4$.
\end{proof}

The proof throws away very little when $\mu$ is spread out: Cauchy--Schwarz is nearly an equality when $|\mu-U|$ is nearly constant. Its strength is that it turns a sum over the $n!$ permutations into a sum over the much fewer irreducible representations, each weighted by
explicit numbers. For the lazy walk on a cycle, it is exactly the computation of Tome~3, Theorem~8.17.

\section{Random transpositions}

The shuffle of Tome~3 picks a card $L$ with the left hand and a card $R$ with the right, independently and uniformly, and swaps them. So the step $\tau$ has law
\[
Q(e)=\frac1n,\qquad Q(\tau)=\frac2{n^2}\ \text{for each transposition }\tau,
\]
and $Q=0$ elsewhere. (Whether we swap cards or positions, that is, multiply on the left or on the right, makes no difference, since $Q$ is constant on conjugacy classes: Exercise~\ref{ex:leftright}.) For $\lambda\vdash n$, write
\[
r(\lambda)=\frac{\chi^\lambda(\tau)}{f^\lambda}=\frac1{\binom n2}\sum_{(i,j)\in\lambda}(j-i),\qquad s(\lambda)=\frac1n+\frac{n-1}n\,r(\lambda).
\]

\begin{proposition}\label{prop:eigenvalues}
$\hat Q(S^\lambda)=s(\lambda)I$. Consequently, the probability that the deck is in the arrangement $\sigma$ after $t$ swaps from the ordered deck is
\[
Q^{*t}(\sigma)=\frac1{n!}\sum_{\lambda\vdash n}f^\lambda\chi^\lambda(\sigma)\,s(\lambda)^t,
\]
and
\[
d(t)^2\le\frac14\sum_{\lambda\ne(n)}(f^\lambda)^2s(\lambda)^{2t}.
\]
\end{proposition}

\begin{proof}
By Proposition~\ref{prop:classfourier}, $\hat Q(S^\lambda)$ is the scalar $\frac1{f^\lambda}\bigl(\frac1nf^\lambda+\binom n2\frac2{n^2}\chi^\lambda(\tau)\bigr)=\frac1n+\frac{n-1}nr(\lambda)$. Then $\widehat{Q^{*t}}(S^\lambda)=s(\lambda)^tI$, and Fourier inversion gives the formula, as characters of
$\mathfrak S_n$ are real (each $\sigma$ is conjugate to $\sigma^{-1}$). The last claim is Lemma~\ref{lem:upperbound}, with $\operatorname{tr}(s^tI\cdot s^tI)=f^\lambda s^{2t}$ and $d_\rho=f^\lambda$.
\end{proof}

So the whole problem is to bound a sum over Young diagrams. The eigenvalue $s(\lambda)$ is close to $1$ only when $\lambda$ has a very long first row: for $\lambda=(n-1,1)$, $r=\frac{n-3}{n-1}$ and $s=1-\frac2n$, and this term alone is $(n-1)^2(1-\frac2n)^{2t}\approx n^2e^{-4t/n}$,
which is small exactly when $t$ exceeds $\frac12n\ln n$ by a multiple of $n$. The work is to show that the other $p(n)-2$ diagrams do not spoil this. Two lemmas do it: one bounds $s(\lambda)$, the other counts.

\begin{lemma}\label{lem:ratiobounds}
Let $\lambda\vdash n$, let $m=\max(\lambda_1,\lambda'_1)$ be the length of its longest row or column, and $k=n-m$. Then
\[
|s(\lambda)|\le\frac mn,\qquad\text{and}\qquad|s(\lambda)|\le\max\Bigl(1-\frac{2k(n-k)}{n^2},\ \frac{k+1}n\Bigr).
\]
\end{lemma}

\begin{proof}
\emph{Crude bounds.} Since $\sum_i\lambda_i(\lambda_i-1)\le(\lambda_1-1)\sum_i\lambda_i=n(\lambda_1-1)$, and similarly for columns, Theorem~\ref{thm:transposition} gives
\[
-\frac{\lambda'_1-1}{n-1}\le r(\lambda)\le\frac{\lambda_1-1}{n-1},\qquad\text{so}\qquad\frac{2-\lambda'_1}n\le s(\lambda)\le\frac{\lambda_1}n.
\]
Both bounds are at most $\frac mn$ in absolute value, which is the first claim.

\emph{Long first row.} Suppose $\lambda_1=n-k$. The other rows hold $k$ cells, so $\sum_{i\ge2}\lambda_i(\lambda_i-1)\le k(k-1)$ (a sum of squares with a fixed sum is largest when concentrated). Dropping the columns,
\[
r(\lambda)\le\frac{(n-k)(n-k-1)+k(k-1)}{n(n-1)}=1-\frac{2k(n-k)}{n(n-1)},\qquad\text{so}\qquad s(\lambda)\le1-\frac{2k(n-k)}{n^2}.
\]
Also $\lambda'_1\le k+1$ (the first column has at most $k$ cells below row $1$), so the crude bound gives $s(\lambda)\ge\frac{1-k}n$. This proves the second claim when $m=\lambda_1$.

\emph{Long first column.} If $m=\lambda'_1$, apply the previous case to $\lambda'$, whose first row is $n-k$. Since contents change sign under conjugation (Corollary~\ref{cor:conjugate}), $r(\lambda)=-r(\lambda')$, and $s(\lambda)=\frac2n-s(\lambda')$. So
\[
-1+\frac2n+\frac{2k(n-k)}{n^2}\le s(\lambda)\le\frac2n+\frac{k-1}n=\frac{k+1}n,
\]
which proves the second claim in this case too.
\end{proof}

\begin{lemma}\label{lem:countshapes}
For $0\le k\le n$, the sum of $(f^\lambda)^2$ over the $\lambda\vdash n$ with $\lambda_1=n-k$ is at most $\binom nk^2k!\le n^{2k}/k!$. The sum over the $\lambda$ with $\max(\lambda_1,\lambda'_1)=n-k$ is at most twice this.
\end{lemma}

\begin{proof}
A standard tableau of shape $\lambda$, with $\lambda_1=n-k$, is determined by the set of $k$ entries below the first row, and by the way they fill the diagram $\mu=(\lambda_2,\lambda_3,\dots)\vdash k$, which, relabelled increasingly by $1,\dots,k$, is a standard tableau of shape $\mu$.
So $f^\lambda\le\binom nkf^\mu$, and distinct $\lambda$ give distinct $\mu$. By Theorem~\ref{thm:rs}, $\sum_{\mu\vdash k}(f^\mu)^2=k!$, which gives the first bound, and $\binom nk\le n^k/k!$ gives the second. Conjugation preserves $f^\lambda$ (Corollary~\ref{cor:conjugate}), and the shapes with
$\lambda'_1=n-k$ are the conjugates of those with $\lambda_1=n-k$.
\end{proof}

\begin{theorem}[Diaconis and Shahshahani, 1981]\label{thm:ds}
For random transpositions on $n\ge10^4$ cards, if $c\ge\frac32$ and $t\ge\frac12n\ln n+cn$, then
\[
d(t)\le3e^{-4c/3}.
\]
\end{theorem}

\begin{proof}
Let $S=\sum_{\lambda\ne(n)}(f^\lambda)^2|s(\lambda)|^{2t}$; by Proposition~\ref{prop:eigenvalues}, $d(t)\le\frac12\sqrt S$. As $|s(\lambda)|\le1$ and $2t\ge N=n\ln n+2cn$, we may replace the exponent $2t$ by $N$. Group the diagrams by $m=\max(\lambda_1,\lambda'_1)=n-k$, and let $x=\frac kn$.

\emph{The sign representation} ($k=0$, $\lambda=(1^n)$): $s=-1+\frac2n$, and $(1-\frac2n)^N\le e^{-2N/n}=n^{-2}e^{-4c}$.

\emph{Long rows or columns: $1\le k\le\frac n3$.} On $[0,\frac13]$, $1-2x(1-x)$ decreases from $1$ to $\frac59$, and $\frac{k+1}n\le\frac13+\frac1n\le\frac59$; so Lemma~\ref{lem:ratiobounds} gives $|s(\lambda)|\le1-2x(1-x)\le e^{-2x(1-x)}$. With Lemma~\ref{lem:countshapes}, these $\lambda$ contribute at most
\[
\frac{2n^{2k}}{k!}\,e^{-2x(1-x)(n\ln n+2cn)}=\frac{2n^{2k}}{k!}\,n^{-2k(1-x)}e^{-4ck(1-x)}\le2\,\frac{n^{2k^2/n}}{k!}\,e^{-8ck/3},
\]
using $1-x\ge\frac23$. Now $k!\ge(k/e)^k$, so $\frac{n^{2k^2/n}}{k!}\le h_k^k$ with $\ln h_k=1+\frac{2k}n\ln n-\ln k$. This is a convex function of $k$, so on $[1,\frac n3]$ it is at most its value at an endpoint: $1+\frac{2\ln n}n\le1+\frac2e$ at $k=1$, and
$1+\ln3-\frac13\ln n\le2.1$ at $k=\frac n3$. So $h_k\le e^{2.1}<8.2$, and with $q=8.2e^{-8c/3}\le8.2e^{-4}<0.151$, the contribution of all these $k$ is at most
\[
2\sum_{k\ge1}q^k=\frac{2q}{1-q}\le2.36\,q<19.4\,e^{-8c/3}.
\]

\emph{Middle shapes: $\frac n3<k\le\frac n2$.} Now $1-2x(1-x)\le\frac59$ and $\frac{k+1}n\le\frac12+\frac1n$, so $|s(\lambda)|\le\sigma_x=\max\bigl(1-2x(1-x),x+\frac1n\bigr)\le0.56$, and $\ln\sigma_x<-0.57$. Using $k!\ge(k/e)^k$ and $\ln k=\ln x+\ln n$, the logarithm of the
contribution of one value of $k$ is at most
\[
\begin{aligned}\ln2+2k\ln n-k\ln k+k+N\ln\sigma_x&=\ln2+n\ln n\,(x+\ln\sigma_x)\\&\quad+nx(1-\ln x)+2cn\ln\sigma_x.\end{aligned}
\]
The function $x+\ln(1-2x+2x^2)$ has derivative $\frac{2x^2+2x-1}{1-2x+2x^2}$, so on $[\frac13,\frac12]$ it is largest at an endpoint, where it is $-0.254$ or $-0.193$; and $x+\ln(x+\frac1n)\le\frac12+\ln(0.5001)<-0.19$. So $x+\ln\sigma_x<-0.19$. Also
$x(1-\ln x)\le0.85$ on this interval. Since $\ln n\ge9.2$, the bound is at most $\ln2-0.89n-1.14cn$, and the $n/6$ or so values of $k$ contribute in total less than $n\,e^{-0.89n}e^{-8c/3}$, which is negligible.

\emph{Short rows and columns: $k>\frac n2$}, that is, $m<\frac n2$. Let $y=\frac mn<\frac12$; then $|s(\lambda)|\le y$ by Lemma~\ref{lem:ratiobounds}. By Lemma~\ref{lem:countshapes}, the weight is at most $2\binom nk^2k!=2n!\binom nk/m!\le2^{n+1}n!/m!$. With $n!\le n^n$ and
$m!\ge(m/e)^m$, the logarithm of the contribution of one value of $m$ is at most
\[
\begin{aligned}(n+1)\ln2+n\ln n-m\ln m+m+N\ln y&=(n+1)\ln2+n\ln n\,(1-y+\ln y)\\&\quad+ny(1-\ln y)+2cn\ln y.\end{aligned}
\]
On $(0,\frac12]$, $1-y+\ln y$ is increasing, hence $\le\frac12-\ln2<-0.19$, $y(1-\ln y)$ is increasing, hence $\le0.85$, and $\ln y\le-\ln2$. With $\ln n\ge9.2$ this is at most $\ln2-0.20n-1.38cn$, and the fewer than $n$ values of $m$ contribute
less than $2n\,e^{-0.2n}e^{-8c/3}$, again negligible for $n\ge10^4$.

\emph{Total.} Adding up,
\[
S\le n^{-2}e^{-4c}+19.4e^{-8c/3}+(\text{negligible})\,e^{-8c/3}\le19.5\,e^{-8c/3},
\]
and $d(t)\le\frac12\sqrt{19.5}\,e^{-4c/3}<3e^{-4c/3}$.
\end{proof}

The threshold $n\ge10^4$ and the constants come from our crude treatment of the middle and short shapes, and are far from optimal. Diaconis and Shahshahani proved $d(\frac12n\ln n+cn)\le ae^{-2c}$ for all $n$ and $c>0$, with a universal constant $a$; the exponent $2c$
is the right one, as the term of $(n-1,1)$ shows. Numerically, for $n=52$ the sum of Proposition~\ref{prop:eigenvalues} already gives $d(155)\le0.06$, with $155\approx\frac12n\ln n+n$ (Figure~\ref{fig:cutoff}).

\begin{corollary}[Cutoff for random transpositions]\label{cor:cutoff}
For $n\ge10^4$,
\[
d\bigl(\tfrac12n\ln n-cn\bigr)\ge1-6e^{1-2c}\quad(c>1),\qquad d\bigl(\tfrac12n\ln n+cn\bigr)\le3e^{-4c/3}\quad(c\ge\tfrac32).
\]
In particular $\frac12n\ln n-2n\le t_{\mathrm{mix}}\le\frac12n\ln n+2n$, and for every $\varepsilon\in(0,1)$, $t_{\mathrm{mix}}(\varepsilon)=\frac12n\ln n+O_\varepsilon(n)$.
\end{corollary}

\begin{proof}
The lower bound is Tome~3, Proposition~8.21, and the upper bound is Theorem~\ref{thm:ds}; at $c=2$, $3e^{-8/3}<0.21<\frac14$. For $\varepsilon\in(0,1)$, choose $c$ with $3e^{-4c/3}\le\varepsilon$ and $c'$ with $1-6e^{1-2c'}>\varepsilon$: then
$\frac12n\ln n-c'n\le t_{\mathrm{mix}}(\varepsilon)\le\frac12n\ln n+cn$.
\end{proof}

So the mixing time is $\frac12n\ln n$, with a window of order $n$, much smaller than $n\ln n$: rescaled by $\frac12n\ln n$, the distance $d$ drops from $1$ to $0$ at time $1$, like a step function. This is the cutoff.

\begin{figure}[ht]
\centering
\fig{ch10-cutoff}
\caption{The distance $d(t)$ to uniform for random transpositions, against $c=(t-\frac12n\ln n)/n$. Solid: the exact value for $n=10$ and $n=20$, computed from the formula of Proposition~\ref{prop:eigenvalues} and the characters of $\mathfrak S_n$. Dashed: the upper bound
$\frac12\sqrt S$ for $n=52$. Dotted: the lower bound $1-6e^{1+2c}$ of Tome~3. In the scale $c$, the curves for different $n$ almost coincide, and the transition from $1$ to $0$ takes place in a window of a few units, that is, a few multiples of $n$ swaps.}
\label{fig:cutoff}
\end{figure}

\begin{think}
Why $\frac12n\ln n$? After $t$ swaps, $2t$ cards have been touched, counted with repetition, and a card never touched is a sure sign of a non-random deck. By the coupon collector of Tome~3, all cards have been touched after about $n\ln n$ touches, that is, $\frac12n\ln n$ swaps.
The lower bound is this argument; the content of the theorem is that nothing else delays mixing. Can you see why the time to touch every card cannot be an upper bound in itself?
\end{think}

\begin{beauty}
Cutoff has since been proved for many shuffles and walks: riffle shuffles (Bayer and Diaconis, 1992, at $\frac32\log_2n$ shuffles; for $52$ cards the distance is still about $0.61$ after six shuffles and $0.33$ after seven, which is where ``seven shuffles'' comes from), random walks on hypercubes, many random walks on groups generated by conjugacy classes, and random walks on random
regular graphs. Yet there is no general theory saying which chains have a cutoff. In 2004 Yuval Peres proposed a simple criterion: the mixing time should be much larger than the inverse of the spectral gap. It is necessary, and it is sufficient for some large classes of
chains, but not in general. For random transpositions, the gap is $\frac2n$ and the mixing time $\frac12n\ln n$; their product $\ln n$ tends to infinity, as the criterion asks.
\end{beauty}

\section{Exercises}

\begin{exercise}[$\star$]\label{ex:leftright}
Show that if $\mu$ is constant on conjugacy classes, then $\mu*\nu=\nu*\mu$ for every $\nu$. Deduce that swapping two random \emph{cards} and swapping the cards in two random \emph{positions} produce the same law of the deck at every time.
\end{exercise}

\begin{exercise}[$\star$]\label{ex:threecards}
For $n=3$, compute $s(\lambda)$ for the three partitions of $3$, and deduce the law of the deck after $t$ swaps and the exact formula $d(t)=\frac{11-(-1)^t}{12}\,3^{-t}$.
\end{exercise}

\begin{exercise}[$\star$]\label{ex:spectralonly}
For $n\ge7$, show that the largest $|s(\lambda)|$, $\lambda\ne(n)$, is $1-\frac2n$, attained at $(n-1,1)$ and $(1^n)$. How many swaps does the spectral bound of Tome~3 (Proposition~8.16) guarantee? Compare with Theorem~\ref{thm:ds}.
\end{exercise}

\begin{exercise}[$\star\star$]\label{ex:fixedpoints}
Show that $\mathbb E[\chi^\lambda(X_t)]=f^\lambda s(\lambda)^t$ for the deck $X_t$ after $t$ swaps from the ordered deck. Deduce that the expected number of cards in their original positions is $1+(n-1)(1-\frac2n)^t$, and comment on the value at $t=\frac12n\ln n+cn$.
\end{exercise}

\begin{exercise}[$\star\star$]\label{ex:lazyhalf}
In the \emph{lazy} version of the shuffle, the deck stays put with probability $\frac12$, and otherwise a uniformly random transposition is applied. Find the eigenvalues $s_{1/2}(\lambda)$, show that the expected number of cards in their original positions after $t$ steps is
$1+(n-1)\bigl(1-\frac1{n-1}\bigr)^t$, and that it is at least $1+e^c$ for $t\le(n-2)(\ln(n-1)-c)$. Why does the lazy walk need about twice as many steps as the shuffle of Theorem~\ref{thm:ds}?
\end{exercise}

\begin{exercise}[$\star\star\star$]\label{ex:orderedpairs}
Let $M$ be the permutation representation of $\mathfrak S_n$ on ordered pairs $(a,b)$ of distinct points, $n\ge4$. (a) Show that
\[
M\cong S^{(n)}\oplus2S^{(n-1,1)}\oplus S^{(n-2,2)}\oplus S^{(n-2,1,1)},
\]
using the remark after Theorem~\ref{thm:specht}, Exercise~\ref{ex:youngmult}, and the fact that for permutation representations on sets $A$ and $B$, $\langle\chi_A,\chi_B\rangle$ is the number of orbits of $\mathfrak S_n$ on $A\times B$. (b) Compute $s(n-2,2)$ and $s(n-2,1,1)$, and
show that the number $F$ of cards in their original positions satisfies
\[
\begin{aligned}\mathbb E\bigl[F(X_t)(F(X_t)-1)\bigr]&=1+2(n-1)\Bigl(1-\frac2n\Bigr)^t\\&\quad+\frac{n(n-3)}2\Bigl(1-\frac2n\Bigr)^{2t}+\frac{(n-1)(n-2)}2\Bigl(1-\frac4n\Bigr)^t.\end{aligned}
\]
\end{exercise}

\section{Hints}

\noindent\textbf{\ref{ex:leftright}.} In $\sum_h\mu(h)\nu(h^{-1}g)$, put $x=h^{-1}g$ and conjugate.

\noindent\textbf{\ref{ex:threecards}.} The contents of $(2,1)$ add up to $0$. Use the formula of Proposition~\ref{prop:eigenvalues} on the three classes.

\noindent\textbf{\ref{ex:spectralonly}.} Use Lemma~\ref{lem:ratiobounds} for the first part, and $\pi_{\min}=1/n!$.

\noindent\textbf{\ref{ex:fixedpoints}.} $\mathbb E[\chi^\lambda(X_t)]=\sum_\sigma Q^{*t}(\sigma)\chi^\lambda(\sigma)$ is the trace of a Fourier transform; and $F=\chi^{(n)}+\chi^{(n-1,1)}$ by Exercise~\ref{ex:standardrep}.

\noindent\textbf{\ref{ex:lazyhalf}.} $s_{1/2}(\lambda)=\frac12+\frac12r(\lambda)$, and $\ln(1-u)\ge-\frac u{1-u}$.

\noindent\textbf{\ref{ex:orderedpairs}.} $M=M^{(n-2,1,1)}$. Count orbits on pairs of ordered pairs (there are $7$) and on (ordered pair, point) (there are $3$). Then $\chi_M(\sigma)=F(\sigma)(F(\sigma)-1)$.

\section{Solutions}

\begin{solution}{ex:leftright}
\idea Conjugation moves the factor $\mu$ from one side to the other.
\writeup Put $x=h^{-1}g$, so $h=gx^{-1}$ and $(\mu*\nu)(g)=\sum_x\mu(gx^{-1})\nu(x)$. Conjugating by $x$, $\mu(gx^{-1})=\mu(x^{-1}(gx^{-1})x)=\mu(x^{-1}g)$, so
\[
(\mu*\nu)(g)=\sum_x\nu(x)\mu(x^{-1}g)=(\nu*\mu)(g).
\] Swapping two cards is $X_{t}=\tau_tX_{t-1}$, with law $Q*\mathcal L(X_{t-1})$;
swapping two positions is $X_t=X_{t-1}\tau_t$, with law $\mathcal L(X_{t-1})*Q$. Since $Q$ is constant on conjugacy classes, the two recursions coincide, and by induction the laws agree at all times.
\end{solution}

\begin{solution}{ex:threecards}
\idea Three eigenvalues, three classes.
\writeup The contents of $(3)$, $(2,1)$, $(1,1,1)$ add up to $3,0,-3$, so $r=1,0,-1$ and $s=1,\frac13,-\frac13$, with $f=1,2,1$. With $a=3^{-t}$ and $b=(-\frac13)^t$, and the characters $(1,1,1)$, $(2,0,-1)$, $(1,-1,1)$ on the classes of $e$, transpositions and $3$-cycles, Proposition~\ref{prop:eigenvalues} gives
\[
Q^{*t}(e)=\frac{1+4a+b}6,\qquad Q^{*t}(\tau)=\frac{1-b}6,\qquad Q^{*t}(\text{3-cycle})=\frac{1-2a+b}6.
\]
(They add up to $1$ over the $1+3+2$ elements.) As $|b|=a$, $4a+b>0$ and $-2a+b<0$, so
\[
d(t)=\frac12\cdot\frac{(4a+b)+3a+2(2a-b)}6=\frac{11a-b}{12}=\frac{11-(-1)^t}{12}\,3^{-t}.
\]
Check: $d(0)=\frac56$, the distance from a point mass to uniform on $6$ points, and $d(1)=\frac13$.
\end{solution}

\begin{solution}{ex:spectralonly}
\idea Only the largest eigenvalue enters the spectral bound, and it pays $\ln n!$ instead of $\ln n$.
\writeup Let $k=n-\max(\lambda_1,\lambda'_1)$ as in Lemma~\ref{lem:ratiobounds}. If $2\le k\le\frac n2$, then $k(n-k)\ge2(n-2)$, so $1-\frac{2k(n-k)}{n^2}\le1-\frac{4(n-2)}{n^2}<1-\frac2n$ (as $n>4$), and $\frac{k+1}n\le\frac12+\frac1n<1-\frac2n$ (as $n>6$). If $k>\frac n2$, then $|s|\le\frac{n-k}n<\frac12$. For $k=1$ the only shapes are $(n-1,1)$ and $(2,1^{n-2})$, with $s=1-\frac2n$ and
$\frac2n-(1-\frac2n)=-1+\frac4n$; and for $k=0$, $(1^n)$ has $s=-1+\frac2n$. So $\lambda_*=1-\frac2n$. The spectral bound gives $d(t)\le\frac{n!}2(1-\frac2n)^t\le\frac{n!}2e^{-2t/n}$, which is $\le\frac14$ once $t\ge\frac n2\ln(2\,n!)\approx\frac12n^2\ln n$: a factor $n$ worse than the truth.
The spectral bound uses only the largest eigenvalue but must pay for all $n!$ dimensions; the upper bound lemma uses every eigenvalue with its exact multiplicity $(f^\lambda)^2$.
\end{solution}

\begin{solution}{ex:fixedpoints}
\idea The expectation of a character is the trace of a Fourier coefficient.
\writeup $\mathbb E[\chi^\lambda(X_t)]=\sum_\sigma Q^{*t}(\sigma)\operatorname{tr}\rho_\lambda(\sigma)=\operatorname{tr}\widehat{Q^{*t}}(S^\lambda)=\operatorname{tr}(s(\lambda)^tI)=f^\lambda s(\lambda)^t$. The number of fixed points is the character of $\C^n=M^{(n-1,1)}\cong S^{(n)}\oplus S^{(n-1,1)}$ (Exercise~\ref{ex:standardrep}), so
$\mathbb E[F(X_t)]=1+(n-1)(1-\frac2n)^t$. At $t=\frac12n\ln n+cn$, $(1-\frac2n)^t\approx e^{-2t/n}=\frac{e^{-2c}}n$, and $\mathbb E[F]\approx1+e^{-2c}$: for $c$ large it is close to $1$, the value under the uniform law; for $c$ very negative it is huge, and the deck is visibly not mixed.
\end{solution}

\begin{solution}{ex:lazyhalf}
\idea The same Fourier analysis with a different weight on the identity; half the steps are wasted.
\writeup The step has law $\frac12\mathbf 1_e+\frac1{2\binom n2}\sum_\tau\mathbf 1_\tau$, with Fourier coefficient $\frac12+\frac12r(\lambda)$ on $S^\lambda$, by Proposition~\ref{prop:classfourier}. For $(n-1,1)$, $r=\frac{n-3}{n-1}$ and $s_{1/2}=\frac{n-2}{n-1}=1-\frac1{n-1}$. As in Exercise~\ref{ex:fixedpoints},
$\mathbb E[F(X_t)]=1+(n-1)(1-\frac1{n-1})^t$. With $u=\frac1{n-1}$, $\ln(1-u)\ge-\frac u{1-u}=-\frac1{n-2}$ (the inequality holds since $\ln(1-u)+\frac u{1-u}$ vanishes at $0$ and has derivative $\frac u{(1-u)^2}\ge0$), so
$\mathbb E[F(X_t)]-1\ge(n-1)e^{-t/(n-2)}\ge e^c$ when $t\le(n-2)(\ln(n-1)-c)$. So the expected number of fixed points is still far from its uniform value $1$, and the upper bound lemma, whose term for $(n-1,1)$ is $(n-1)^2s_{1/2}^{2t}=(\mathbb E[F]-1)^2\ge e^{2c}$, cannot prove mixing before about $n\ln n$ steps. The lazy walk performs a transposition only half the time, and a uniform transposition
corresponds to $\frac{n-1}n$ of a step of the original walk (which wastes a step with probability $\frac1n$); so $t$ lazy steps are worth about $\frac t2$ steps of the original, and $n\ln n$ lazy steps are needed instead of $\frac12n\ln n$.
\end{solution}

\begin{solution}{ex:orderedpairs}
\idea Dominance limits the constituents to four; orbit counts fix the multiplicities. Then each constituent contributes $f^\lambda s(\lambda)^t$.
\writeup (a) Tabloids of shape $(n-2,1,1)$ are determined by their second and third rows, an ordered pair of distinct points, so $M=M^{(n-2,1,1)}$. The partitions dominating $(n-2,1,1)$ are $(n)$, $(n-1,1)$, $(n-2,2)$ and $(n-2,1,1)$; so
$M\cong m_1S^{(n)}\oplus m_2S^{(n-1,1)}\oplus m_3S^{(n-2,2)}\oplus m_4S^{(n-2,1,1)}$ with $m_4=1$ by Exercise~\ref{ex:youngmult}. The orbits of $\mathfrak S_n$ on pairs $((a,b),(c,d))$ are classified by the coincidences between $\{a,b\}$ and $\{c,d\}$: $(c,d)=(a,b)$; $(c,d)=(b,a)$; exactly one coincidence,
$c=a$, $c=b$, $d=a$ or $d=b$, the other point new; or no coincidence (possible since $n\ge4$). That is $7$ orbits, so $m_1^2+m_2^2+m_3^2+m_4^2=\langle\chi_M,\chi_M\rangle=7$. The action on $M$'s basis is transitive, so $m_1=\langle\chi_M,1\rangle=1$. The orbits on pairs $((a,b),c)$ are $c=a$, $c=b$, and $c$ new:
$3$ of them, and $\C^n\cong S^{(n)}\oplus S^{(n-1,1)}$, so $m_1+m_2=3$ and $m_2=2$. Then $m_3^2=7-1-4-1=1$.

(b) The contents of $(n-2,2)$ are $0,\dots,n-3$ and $-1,0$, with sum $\binom{n-2}2-1$; so $r=\frac{(n-2)(n-3)-2}{n(n-1)}=\frac{(n-1)(n-4)}{n(n-1)}=\frac{n-4}n$ and $s=\frac1n+\frac{(n-1)(n-4)}{n^2}=\frac{(n-2)^2}{n^2}=(1-\frac2n)^2$. Those of $(n-2,1,1)$ are $0,\dots,n-3$ and $-1,-2$, with sum $\binom{n-2}2-3$;
so $r=\frac{n^2-5n}{n(n-1)}=\frac{n-5}{n-1}$ and $s=\frac1n+\frac{n-5}n=1-\frac4n$. The hook length formula gives $f^{(n-2,2)}=\frac{n(n-3)}2$ and $f^{(n-2,1,1)}=\frac{(n-1)(n-2)}2$. An ordered pair of distinct points is fixed by $\sigma$ if and only if both points are fixed, so $\chi_M(\sigma)=F(\sigma)(F(\sigma)-1)$, and
by (a) and Exercise~\ref{ex:fixedpoints},
\[
\mathbb E[F(X_t)(F(X_t)-1)]=\sum_\lambda m_\lambda f^\lambda s(\lambda)^t,
\]
which is the stated formula. At $t=0$ it gives $1+2(n-1)+\frac{n(n-3)}2+\frac{(n-1)(n-2)}2=n(n-1)$, as it should. With Exercise~\ref{ex:fixedpoints}, this gives the variance of $F(X_t)$ exactly, and a second-moment proof of the lower bound of Corollary~\ref{cor:cutoff}.
\end{solution}

\section*{Sources and further reading}
\addcontentsline{toc}{section}{Sources and further reading}

P.~Diaconis and M.~Shahshahani, ``Generating a random permutation with random transpositions'', \emph{Z.\ Wahrscheinlichkeitstheorie verw.\ Gebiete} 57 (1981), 159--179; D.~Aldous and P.~Diaconis, ``Shuffling cards and stopping times'', \emph{Amer.\ Math.\ Monthly} 93 (1986), 333--348;
D.~Bayer and P.~Diaconis, ``Trailing the dovetail shuffle to its lair'', \emph{Ann.\ Appl.\ Probab.} 2 (1992), 294--313. The biographical facts and quotations about Diaconis, Vernon, Gardner and Mosteller are from the MacTutor biography of Persi Diaconis (J.~J.~O'Connor and E.~F.~Robertson).

For more: P.~Diaconis, \emph{Group Representations in Probability and Statistics} (Institute of Mathematical Statistics, 1988), Chapters~3 and~4, the source of this chapter; D.~A.~Levin and Y.~Peres, \emph{Markov Chains and Mixing Times} (American Mathematical Society, 2nd edition 2017), Chapters~8, 18 and~23, on
cutoff and Peres' criterion; L.~Saloff-Coste, ``Random walks on finite groups'', in \emph{Probability on Discrete Structures}, Encyclopaedia Math.\ Sci.\ 110 (Springer, 2004), 263--346.
